\subsection{محدودیت‌ها و تحلیل شکست}
با وجود بازدهی و سرعت بالا، روش تجزیه کرونکر دارای محدودیت‌های تئوریک و عملیاتی زیر است:
\begin{itemize}
    \item \textbf{انعطاف‌ناپذیری ابعادی در ابعاد اول:} کیفیت فاکتورگیری کرونکر شدیداً وابسته به تجزیه‌پذیری ابعاد $d_{\text{in}}$ و $d_{\text{out}}$ به فاکتورهای غیربدیهی نزدیک به جذر است. در لایه‌هایی با ابعاد عدد اول یا فاکتورهای نامتعادل، فشردگی کاهش یافته یا لایه ناگزیر به مسیر رزرو \lr{SVD} هدایت می‌شود.
    \item \textbf{افت دقت در دیکودرهای خودبازگشتی طویل:} در مدل‌های زبانی مانند \lr{GPT-2} و \lr{Qwen}، ماهیت ساختار کرونکر قیدهای صلب بلوکی روی درایه‌های وزن ایجاد می‌کند؛ در نتیجه در غیاب تیونینگ کافی، خطای تجمعی در گام‌های توالی طولانی سبب افت کیفیت پیش‌بینی لایه \lr{LM Head} می‌شود.
    \item \textbf{سربار جابجایی تانسورها (\lr{Permutation and Reshape Overheads}):} در پیاده‌سازی گذر پیشرو بدون ماتریسی‌سازی صریح، نیاز به عملیات ترانهاده و تغییر شکل مکرر داده‌ها در حافظه کش وجود دارد که در توالی‌های کوتاه می‌تواند مانع از رسیدن به حداکثر شتاب تئوریک سخت‌افزاری گردد.
\end{itemize}
